% Source presentations for register row N-c-004.

\subsection{Four descriptive-complexity presentations}
\label{sec:p09-presentation-n-c-004}

\subsubsection{Fagin's generalized spectra}

Fagin begins with two disjoint finite relational similarity types \(S\) and
\(T\).  If \(\theta\) is a first-order sentence over \(S\cup T\), then
\(\operatorname{Mod}_{S}\exists T\theta\) is the class of finite
\(S\)-reducts which admit, on the same universe, interpretations of the
relations in \(T\) making \(\theta\) true.  Such a class is an
\((S,T)\)-spectrum; it is an \(S\)-spectrum when some finite \(T\) supplies
that representation.  The traditional numerical spectrum is the case
\(S=\varnothing\) \cite[pp.~48--49]{Fagin1974}.  Here ``spectrum'' names a
class of finite reducts, not a resource bound or a set of models of arbitrary
cardinality.

For nonempty \(S\), a finite \(S\)-structure \(\mathcal A\) on
\(\{1,\ldots,n\}\) is encoded by \(e(\mathcal A)\): the binary expansion of
\(n\), a separator, and the lexicographically ordered characteristic words of
its relations.  Theorem~6 says that an isomorphism-closed class of finite
\(S\)-structures is an \(S\)-spectrum exactly when the language of these
encodings is accepted by a nondeterministic polynomial-time machine.  When
\(S=\varnothing\), the source still uses the binary representation of size,
and its class \(\mathsf{NP}_1=\bigcup_k\mathsf{NTIME}(2^{k\ell})\)
measures linear-exponential time in that bit length. Polynomial time after
unary recoding is a separate consequence, not the source's input encoding
\cite[pp.~51--54]{Fagin1974}. The complete source-specific statement and
both constructions are \cref{thm:p09-fagin74-theorem6}. The modern slogan
``existential second-order logic captures \(\mathsf{NP}\)'' is therefore an
equivalent summary of an encoding- and isomorphism-sensitive finite-structure
theorem, not the source's whole statement.

Both directions are constructed.  From an existential second-order sentence,
the machine guesses the finite relation tables and checks the fixed first-order
matrix on all assignments in polynomial time.  Conversely, after normalizing a
nondeterministic multitape machine, existentially quantified relations encode
an auxiliary order, time and tape coordinates, nondeterministic choices, and a
computation tableau.  The first-order matrix enforces the initial row, every
local transition, and acceptance \cite[pp.~53--58]{Fagin1974}. The auxiliary
order is guessed as part of the expansion; the input structure is not assumed
to carry a built-in order.  Even cardinality, expressed by the existence of a
fixed-point-free involution, is one of the source's elementary spectrum
examples.  The public handbook identity for the result remains
Theorem~\ref{thm:p03-fagin}.

\subsubsection{Vardi's query-complexity matrix}

Vardi takes a finite database to be
\(B=(D,R_1,\ldots,R_k)\), with \(D\subseteq\mathbb N\), and treats a query as
an isomorphism-consistent partial function on databases. Its tuple-recognition
graph records answer membership, but does not distinguish an undefined query
from a defined empty answer. Definable-language capture therefore concerns
total consistent queries, or queries identified by those graphs, rather than
unqualified equality of partial-function domains. The source separately varies
the database, the expression, or
both.  These give \emph{data complexity}, \emph{expression complexity}, and
\emph{combined complexity}, respectively.  They are three input-size
conventions for one query-evaluation problem and cannot be exchanged without
changing the claim \cite[pp.~137--139]{Vardi1982}.

The source notation \(F\), \(TF\), \(YF\), and \(SF\) names four query
languages: first-order relational calculus, first order with transitive
closure, first order with an outer positive least-fixed-point operator, and existential
second-order formulas.  These abbreviations are retained here only to report
Vardi's source results; they are not new handbook-wide language names.
Theorem~5 makes the data complexity of \(YF\) complete for polynomial time and
its expression complexity complete for exponential time.  For an \(r\)-ary
positive operator on a database of size \(n\), the upper bound follows by at
most \(n^r\) fixed-point stages. The data-complexity lower-bound sketch
uses Path System Accessibility. The separate expression-complexity
lower-bound sketch describes a machine tableau on a fixed database; these
are not one lower-bound argument
\cite[pp.~141--142]{Vardi1982}.

The ordered-capture result appears as a qualified remark rather than as the
later textbook theorem in modern syntax.  Vardi permits a slight extension
\(YF^+\) in which existential quantification may range over conjunctions of
equality formulas with \(YF\) expressions.  If a linear order is already in
the database, write the source's weak-order relation as \(\leq\).
At total consistent-query scope the capture statement is
\[
 C(YF^+,\leq)=QPTIME_{\rm tot}(\leq),
\]
where the left side is the class of ordered queries definable in that source
language and the right side is the class of total polynomial-time queries on ordered
databases.  The deterministic machine tableau supplies the proof sketch
\cite[remark~2]{Vardi1982}.  The paper is an extended abstract: the iteration
bound and tableau idea are given, but a line-by-line normal-form simulation is
not.  Its `Even' query shows that unordered positive fixed point does not
capture every polynomial-time query, so the order hypothesis is load-bearing.

\subsubsection{Immerman's 1986 fixed-point theorem}

Immerman writes a finite relational database as
\(B=\langle U,R_1,\ldots,R_k,c_1,\ldots,c_s\rangle\).  His language \(L\) is
first-order relational calculus.  The language \(L_\mu\) closes \(L\) under
least fixed points of formulas positive in the relation being defined; the
positivity requirement makes the induced operator monotone.  The main setup
assumes at least one constant \(c_0\).  Without a constant, the two capture
theorems require one final quantifier to close the remaining free variables
\cite[pp.~87--89]{Immerman1986}.  These language symbols are source-specific
query languages and are used below only to state the source's theorem.

Fact~2.1 proves that every \(L_\mu\) query is computable in polynomial time,
because an \(r\)-ary relation can gain a tuple at most \(\lvert U\rvert^r\)
times.  Fact~2.2 uses parity to show that unordered \(L_\mu\) is strictly weaker
than the class of polynomial-time queries.  With a distinguished total order,
Theorem~1 identifies polynomial-time queries with \(L_\mu\)-definable queries
\cite[pp.~89--93]{Immerman1986}.  The source's displayed equation also uses its
pre-existing notation \(\Sigma_\omega(\leq)\); that symbol is not promoted to
handbook notation because the equality between ordered polynomial time and
\(L_\mu(\leq)\) is the part used here.

The hard inclusion is proved in detail.  Identify the ordered universe with
\(0,\ldots,n-1\), use \(k\)-tuples to address \(n^k\) time and tape positions,
and define by one least fixed point the relation describing the deterministic
computation tableau.  First-order formulas define its initial row and its local
successor configurations; a designated state/bit records acceptance.  The
proof treats Boolean queries for simplicity and works through the tableau
construction rather than merely citing it \cite[pp.~90--93]{Immerman1986}.
This is the proof route underlying the public
Theorem~\ref{thm:p04-immerman-vardi}; it is not Vardi's extended-abstract
presentation and it does not remove the order hypothesis.

\subsubsection{Immerman's 1999 existential-second-order presentation}

In Chapter~7 of \emph{Descriptive Complexity}, \(SO\exists\) denotes the
Boolean queries definable by existential second-order sentences.  The book
first gives SAT, CLIQUE, and graph 3-colourability examples, making the guessed
relations concrete.  Proposition~7.6 proves \(SO\exists\subseteq\mathsf{NP}\)
by guessing each finite relation table and model-checking the fixed first-order
matrix.  Theorem~7.8 proves the converse and strengthens the normal form: the
first-order matrix may be chosen universal
\cite[pp.~143--148]{Immerman1999}.

For the converse, existential relations \(C_i\) of arity \(2k\) encode the
machine configurations and a \(k\)-ary relation \(\Delta\) records the
nondeterministic choices.  Universal first-order conjuncts impose the initial
configuration, consistency of those choices, local transitions, and an
accepting configuration.  This is a complete pedagogical tableau proof and is
more explicit than the 1974 paper's original spectrum notation.  It remains a
distinct presentation: its examples, \(SO\exists\) syntax, and universal-matrix
refinement are not retrojected into Fagin's article.

The same registered book locates its ordered fixed-point treatment in
Chapter~4, but the technical text of that chapter was not available in the
verified source copy.  Its theorem number, exact syntax, hypotheses, and proof
are therefore not supplied here from the 1986 article.  The 1999 presentation
above is complete for the Chapter~7 existential-second-order result; it does
not supply the distinct Chapter~4 fixed-point presentation
\cite{Immerman1999}.
